% Источники: exam/materials/моделирование.pdf, раздел 24; https://github.com/HumsterProgrammer/iu9-notes/blob/main/8sem/Моделирование/Моделирование_лекция-05_26-02-26.md.
\section{Внутренняя синхронизация имитационной модели.}

\textbf{Внутренняя синхронизация} компонент модели требуется, когда для каждой компоненты сложной системы, кроме алгоритма, необходимо установить временную координату функционирования. Добавление операторов преобразования временной координаты компоненты модели является внутренней синхронизацией.

\subsection*{Семафоры}

\textbf{Семафоры} -- глобальные переменные, значение которых либо разрешает, либо запрещает активизацию имитатора компоненты в модельном времени.

В случае запрета активизации компоненты модели выстраиваются в очередь на активизацию согласно их приоритетам. Сброс семафора приводит к активизации ожидающих его компонент; установка семафора порождает очередь компонент на активизацию.

\subsection*{Очереди}

Для транзактного способа имитации сами очереди к блокам являются средством внутренней синхронизации. В СМО могут использоваться дисциплины обслуживания:
\begin{itemize}
  \item FIFO -- обслуживание в порядке поступления;
  \item LIFO -- обслуживание по принципу стека;
  \item случайный выбор заявок;
  \item обслуживание по приоритету.
\end{itemize}

При процессном способе, кроме семафоров, необходимы установка операторов взаимодействия процессов друг с другом и выбор дисциплин обслуживания очередей.

\clearpage
